\section{Irreducible polynomials over \texorpdfstring{$\F_2$}{F2}}
\label{sec:field-existence}

\begin{lemma}\label{lem:field-existence}
For every integer $r\ge1$, there is a monic irreducible polynomial of
degree $r$ over $\F_2$.
\end{lemma}

\begin{proof}
First construct an extension field in which $X^{2^r}-X$ splits.  Any
nonconstant polynomial over a field has an irreducible factor, by choosing
a nonconstant divisor of least degree.  The quotient by an irreducible
polynomial is a field: the Euclidean algorithm expresses $1$ as a linear
combination of that polynomial and any nonzero residue representative,
giving the inverse of the latter.  Adjoining a root by this quotient and
factoring out the corresponding linear factor can be repeated.  Each step
reduces the degree still to be split, so finitely many such extensions
suffice.

In the resulting field the roots of $X^{2^r}-X$ are distinct, since its
formal derivative is $-1=1$.  There are exactly $2^r$ roots.  Their set
$F$ contains $0,1$ and is closed under addition and multiplication: in
characteristic two, repeated squaring gives
$(a+b)^{2^r}=a^{2^r}+b^{2^r}$, and the analogous product identity is
immediate.  For a nonzero root $a$, the identity
$(a^{-1})^{2^r}=(a^{2^r})^{-1}=a^{-1}$ gives closure under inversion.
Negatives equal the original elements in characteristic two.  Thus $F$
is a field with $2^r$ elements.  A basis of $F$ over $\F_2$ has exactly
$r$ elements, because a vector space of dimension $v$ over $\F_2$ has
$2^v$ elements.

For $a\in F$, let its monic polynomial of smallest positive degree over
$\F_2$ that vanishes at $a$ have degree $e$.  Such a polynomial exists by
linear dependence in $F$, and it is irreducible: a proper factorization
would give a smaller vanishing factor in the field $F$.  Polynomial
division shows that evaluation identifies $\F_2[a]$ with its quotient
field, with basis $1,a,\ldots,a^{e-1}$.  It therefore has $2^e$ elements.
Regarding $F$ as a vector space over this subfield shows $e\mid r$:
if its dimension is $v$, the products of the two bases give an
$\F_2$-basis of size $ev=r$.

Every element of a field with $2^e$ elements satisfies $a^{2^e}=a$.
For nonzero $a$, multiplication by $a$ permutes the nonzero elements;
taking their product and cancelling gives $a^{2^e-1}=1$.  Zero satisfies
the identity as well.  Hence each element of $F$ of degree $e$ is among
the at most $2^e$ roots of $X^{2^e}-X$.  Here the usual root bound follows
by repeatedly dividing by $X-a$ for distinct roots.  If $e<r$ and
$e\mid r$, then $e\le\floor{r/2}$.  The number of elements of degree
less than $r$ is consequently at most
\[
 \sum_{e=1}^{\floor{r/2}}2^e<2^r,
\]
with the sum empty for $r=1$.  Some element of $F$ has degree $r$, and
its monic minimal polynomial is the required irreducible polynomial.
\end{proof}
